\section{Introduction}

The advancement of intelligent transportation systems (ITS) has accelerated the
deployment of Vehicle-to-Everything (V2X) communications, enabling vehicles to
exchange real-time information for safety and traffic efficiency. Among V2X paradigms,
Vehicle-to-Vehicle (V2V) communication is particularly critical for safety-related
applications such as collision avoidance and emergency braking, especially in
high-speed highway environments where authentication latency directly impacts safety
outcomes.

In direct V2V (sidelink) communication, vehicles exchange safety messages---such as
Basic Safety Messages (BSMs)---directly over the air without relying on cellular
infrastructure or roadside units. Vehicles must authenticate messages from
neighbouring vehicles rapidly and reliably. The IEEE 802.11p standard and its
successor IEEE 802.11bd impose strict latency constraints, typically requiring
end-to-end message processing within 100--300\,ms, of which authentication
constitutes a significant portion.

\subsection{Problem Statement}

Existing V2V authentication mechanisms face several specific limitations in
safety-critical highway environments:

\textbf{Certificate-based schemes} (e.g., ECDSA with X.509 certificates as used in
IEEE 1609.2/SCMS) provide strong security guarantees but introduce non-negligible
overhead. Each safety message requires signature generation, certificate attachment,
and full signature verification at the receiver. In dense traffic scenarios where a
vehicle may receive hundreds of BSMs per second, the cumulative verification cost
creates processing bottlenecks. Certificate revocation relies on Certificate
Revocation Lists (CRLs) that grow as $O(n)$ in the number of revoked pseudonyms,
imposing significant storage and distribution overhead on the vehicular
network~\cite{scopelliti2024}.

\textbf{Signature-based and group signature schemes} reduce some certificate
management overhead but often rely on computationally expensive bilinear pairing
operations. Batch verification techniques have been proposed for signature-based
schemes~\cite{zhang2024}, but these approaches still expose pseudonymous identifiers
linkable across messages within a pseudonym's validity period, leaving vehicles
vulnerable to trajectory tracking.

\textbf{Existing ZKP-based vehicular authentication schemes} have emerged as a
promising direction for privacy preservation. However, these approaches have not
adequately addressed two critical aspects: (i)~the optimization of \emph{frequent
authentication transactions}---vehicles broadcast safety messages every 100--300\,ms,
requiring authentication at high frequency, and existing ZKP schemes treat each
authentication as an independent, full-cost operation; and (ii)~\emph{efficient
anonymous credential management}---existing schemes either rely on blockchain
consensus for revocation (introducing latency) or do not provide a lightweight
mechanism for verifying credential validity without revealing
identity~\cite{jiang2024,pathak2024,zheng2024hybrid}.

The specific open challenge we address is: \emph{how to design a V2V authentication
protocol that minimizes per-message authentication cost through proof precomputation
and session-based credential management, while maintaining unlinkability, efficient
revocation, and scalability under dense traffic conditions.}

\subsection{Proposed Solution and Technical Contributions}

We propose \textbf{ULP-V2V-Auth} (Ultra-Low-Latency Privacy-Preserving V2V
Authentication), a ZKP-based authentication framework whose novelty lies not in the
use of ZKP itself, but in three specific design choices that jointly optimize the
authentication pipeline for high-frequency V2V:

\begin{enumerate}
    \item \textbf{Anonymous Session Token (AST) Architecture with Merkle Tree
    Revocation.} Instead of authenticating each message independently using
    long-term credentials, vehicles acquire short-lived Anonymous Session Tokens
    during low-activity periods. The AST acts as a session-level anonymous
    credential, allowing multiple safety messages to be authenticated under a
    single token without re-engaging the Trusted Authority. Revocation is managed
    via a Merkle hash tree, where each vehicle only needs to maintain its own
    Merkle inclusion proof ($O(\log n)$ in the number of active tokens) rather than
    downloading a full CRL.

    \item \textbf{Groth16 zk-SNARK with Proof-Slot Cache.} We decompose ZKP
    generation into an offline precomputation phase and a near-zero online phase.
    Circuit-level analysis shows that 97.3\% of constraints are message-independent;
    leveraging this, complete Groth16 proofs are precomputed offline during AST
    acquisition---each committed to a predicted BSM content---and stored in a proof
    cache. The online phase reduces to a single Poseidon-2 hash (0.44\,ms), achieving
    sub-1\,ms per-message authentication latency on OBU-class hardware.

    \item \textbf{Batch Verification with DCV Fallback for Dense Traffic Scalability.}
    Receiving vehicles aggregate multiple incoming ZKP proofs and verify them in a
    single batched pairing check, exploiting the algebraic structure of Groth16. This
    reduces per-proof verification cost from $3k$ pairings (individual) to
    approximately $k+3$ pairings (batched), directly addressing the scalability
    bottleneck in dense traffic. Our RPi~4 experiments confirm a $\mathbf{4.65\times}$
    measured speedup at $k=30$, exceeding the theoretical $2.73\times$ pairing-count
    ratio. We additionally adapt the divide-and-conquer verifier
    (DCV)~\cite{ferrara2008finding} to the Groth16 setting, bounding adversarial
    batch-injection fallback cost at $O(w \log k)$ sub-batch operations for $w$
    invalid proofs---a $4.2\times$ improvement over naive $O(k)$ individual fallback
    at $k=30$.
\end{enumerate}
